\section{Results I: the corpus and the dataset}
\label{sec:results-data}

\subsection{Search and screening}

The four searches returned 5,731 records, of which 2,141 were duplicates (\cref{fig:prisma}). Of the 3,590 unique
records, 2,493 were excluded on title and abstract, most often because the study did not use an unloading model
(1,048) or reported no muscle outcome (830); 481 of the latter were real unloading campaigns that reported only gait,
motor-unit or cardiovascular outcomes, and were kept as leads for the campaign map. Seventy-four records were
included at title and abstract; 70 full texts were retrieved (30 as machine-readable XML, 40 as PDF) and four were
not, two because they were paywalled and two because they exist only as conference abstracts. One of the latter still
contributes: its abstract prints a complete result, which was extracted with low confidence~\citep{dirks2016}. The
search contributed 42 studies, 572 rows and 29 campaigns to the dataset; 19 included studies publish their muscle
results only as charts without a baseline value and could not be extracted.

A further 1,023 records---the remainder of the \emph{maybe} set and 80 priority records that need a full text
before a decision can be made---were not screened in the time available. They are not exclusions and are not
reported as such. \draftnote{the full-text stage of the PRISMA table (\file{docs/literature-review/prisma\_counts.md})
still reads ``pending'' for the number of studies included, and its count of 11 full-text exclusions carries
reasons for only 5. The diagram shows the counts as recorded; they should be reconciled with the final extraction
before the report is circulated.}

\begin{figure}[htbp]
\centering
\begin{tikzpicture}[node distance=4.5mm and 6mm, every node/.style={font=\sffamily\scriptsize}]
  % main column
  % the side column is stacked first, and each main-column box is aligned with its side boxes
  \node[box=6.0cm] (id) {\textbf{Records identified} (n = 5,731)\\ PubMed 1,412 \textperiodcentered{} Scopus 1,757\\ Web of Science 1,876 \textperiodcentered{} NASA NTRS 686};
  \node[sidebox=6.6cm, right=8mm of id] (dup) {Duplicates removed (n = 2,141)};
  \node[sidebox=6.6cm, below=6mm of dup] (exc1) {\textbf{Excluded} (n = 2,493): not an unloading model 1,048; no muscle
     outcome 830; review or editorial 346; not human 252; disease causes wasting 12; duration too short 5};
  \node[sidebox=6.6cm, below=2mm of exc1] (unscr) {\textbf{Not yet screened} (n = 1,023): rest of the \emph{maybe} set
     and 80 priority records awaiting a full text};
  \node[sidebox=6.6cm, below=6mm of unscr] (notret) {\textbf{Not retrieved} (n = 4): 2 paywalled, 2 conference abstracts only};
  \node[sidebox=6.6cm, below=5mm of notret] (exc2) {\textbf{Excluded at full text} (n = 11; reasons recorded for 5:
     no muscle outcome 2, no full text 3)\\[1pt] \textbf{Included but not extractable} (n = 19): figure-only, no baseline};
  \node[sidebox=6.6cm, below=6mm of exc2] (other) {\textbf{Other sources}\\ Pre-search corpus: 15 documents $\rightarrow$
     9 studies \textperiodcentered{} 165 rows \textperiodcentered{} 9 campaigns (6 not found by the search)\\
     NASA NLSP open data: 7 campaign folders $\rightarrow$ 1 campaign \textperiodcentered{} 5 rows};
  \coordinate (scrlevel) at ($(exc1.north)!0.5!(unscr.south)$);
  \node[box=6.0cm] (scr) at (id |- scrlevel) {\textbf{Records screened} on title and abstract\\ (n = 3,590)};
  \node[box=6.0cm] (ft) at (id |- notret) {\textbf{Full texts sought} (n = 74)};
  \node[box=6.0cm] (ass) at (id |- exc2) {\textbf{Full texts assessed} (n = 64)};
  \node[keybox=6.0cm] (inc) at (id |- other) {\textbf{Studies from the search contributing rows}\\ 42 studies \textperiodcentered{} 572 rows \textperiodcentered{} 29 campaigns};
  \node[keybox=13.4cm, below=8mm of inc.south west, anchor=north west] (ds) {\textbf{Frozen dataset v1.1}: 52 studies \textperiodcentered{} 742 rows \textperiodcentered{} 36 independent campaigns};
  \node[keybox=13.4cm, below=5mm of ds] (sub) {\textbf{Modelling subset A}: 41 studies \textperiodcentered{} 346 rows \textperiodcentered{} 32 campaigns\\
     subset B: 31 studies \textperiodcentered{} 304 rows \textperiodcentered{} 25 campaigns};
  % arrows
  \draw[flow] (id) -- (scr);
  \draw[flow] (scr) -- node[right, font=\sffamily\scriptsize, text=inktwo] {74 included} (ft);
  \draw[flow] (ft) -- (ass);
  \draw[flow] (ass) -- (inc);
  \draw[flow] (id.east) -- (dup.west);
  \draw[flow] (scr.east) -- ++(4mm,0) |- (exc1.west);
  \draw[flow] (scr.east) -- ++(4mm,0) |- (unscr.west);
  \draw[flow] (ft.east) -- (notret.west);
  \draw[flow] (ass.east) -- (exc2.west);
  \draw[flow] (inc.south) -- (inc.south |- ds.north);
  \draw[flow] (other.south) -- (other.south |- ds.north);
  \draw[flow] (ds) -- (sub);
  % stage labels
  \node[stage] at ($(id.west)+(-3mm,0)$) {Identification};
  \node[stage] at ($(scr.west)+(-3mm,-5mm)$) {Screening};
  \node[stage] at ($(ft.west)!0.5!(ass.west)+(-3mm,0)$) {Full text};
  \node[stage] at ($(inc.west)!0.5!(sub.west)+(-3mm,0)$) {Included};
\end{tikzpicture}
\caption[Flow of records from search to modelling subset]{Flow of records from the search to the modelling
subset, in the style of PRISMA 2020~\citep{page2021}. The search was limited to publications from 2013 onwards;
older work enters only through the pre-search corpus. Six of the nine corpus campaigns were not found by the search,
because their papers predate the date limit.}
\label{fig:prisma}
\end{figure}

\subsection{The dataset}

The frozen dataset (\cref{tab:dataset}) holds 742 comparisons against baseline from 52 studies published between
1992 and 2026 and 36 independent campaigns. Most rows are MRI volumes of named muscles in young men during
head-down tilt bed rest, but every other design, modality and population in the corpus is represented and marked.
The dataset is strongly clustered: 15 campaigns contribute five rows or fewer, while the MEDES 90-day campaign
contributes 299 rows (40\%) to the full file. \Cref{fig:campaigns} shows the 32 campaigns of the modelling subset,
their planned lengths and the days on which they scanned; \cref{app:campaigns} lists all 36.

\begin{table}[htbp]
\centering
\caption[Characteristics of the frozen dataset]{Characteristics of \code{dataset\_v1.1}. Counts are rows unless stated.}
\label{tab:dataset}
\small
\begin{tabularx}{\linewidth}{@{}L{4.4cm} Y@{}}
\toprule
Characteristic & Value \\
\midrule
\inputtable{generated/tab_dataset.tex}
\bottomrule
\end{tabularx}
\end{table}

\paragraph{Origin of the rows} 572 rows (42 studies, 29 campaigns) come from the systematic search, 165 (9 studies,
9 campaigns) from the pre-search corpus and 5 (one campaign) from the NASA archive. Every literature row names its
DOI, source file and page, table or figure; the NASA rows name their folder and the two columns they were computed
from.

\paragraph{Quality markers} 674 rows carry high extraction confidence, 49 medium and 19 low; the 19 low-confidence
rows are the ones read off figures. 25 partial rows lack a group size or a dispersion. 547 rows do not state which
leg was imaged. No row has been double-extracted by a second person (\cref{sec:verification}).

\begin{figure}[htbp]
\centering
\includegraphics[width=\linewidth]{fig_campaigns}
\caption[The 32 campaigns of the modelling subset]{The 32 campaigns of the modelling subset, ordered by planned
length. Each bar runs to the campaign's planned length of unloading; each dot is a scan day, with area proportional to
the number of rows measured that day; the number is the campaign's rows in subset A. Most campaigns scanned only at
the end; the long campaigns that scanned repeatedly carry most of the within-campaign information about time.}
\label{fig:campaigns}
\end{figure}

\subsection{The modelling subsets}

\Cref{tab:subsets} walks the dataset down to the two modelling subsets. Restricting to rows measured during
unloading removes 264 recovery rows, and the spaceflight and recovery-only campaigns with them; restricting to the
lower limb removes 53 trunk-muscle rows; and the one-tissue-one-row rule removes 79 composite rows whose components
were reported on the same occasion. Of the 346 rows in subset A, 246 are from control arms (30 campaigns) and 100
from countermeasure arms; 42 are whole-segment measurements; 251 are MRI volumes. Subset B removes the whole-segment
rows and with them seven campaigns that report nothing else.

\begin{table}[htbp]
\centering
\caption[From the frozen dataset to the modelling subsets]{From the frozen dataset to the modelling subsets. The
predicates are declared in \file{framework/config.yaml} and applied in this order.}
\label{tab:subsets}
\small
\begin{tabularx}{\linewidth}{@{}Y rrr@{}}
\toprule
Step & Rows & Studies & Campaigns \\
\midrule
\inputtable{generated/tab_subsets.tex}
\bottomrule
\end{tabularx}
\end{table}

\subsection{A first look at the data}

Before any model, the unadjusted control-arm means already show the pattern the models go on to quantify
(\cref{tab:classmeans}): antigravity extensors lose about twice as much as flexors. By family, the plantar flexors
lose most ($-14.7\%$), followed by the evertors ($-12.2\%$), knee extensors ($-8.6\%$), knee flexors ($-7.8\%$),
dorsiflexors ($-7.5\%$), hip adductors ($-4.1\%$) and hip rotators ($-1.3\%$). These means ignore duration, arm,
modality and campaign, and are shown only as a sanity check on the muscle classification.

\begin{table}[htbp]
\centering
\caption[Unadjusted control-arm means by functional class]{Unadjusted mean percentage change of control-arm rows
by functional class (lower-limb rows during unloading, before the one-tissue-one-row rule).}
\label{tab:classmeans}
\small
\begin{tabular}{@{}lrr@{}}
\toprule
Functional class & Rows & Mean change (\%) \\
\midrule
\inputtable{generated/tab_class_means.tex}
\bottomrule
\end{tabular}
\end{table}
